The doctrine of being
§ 84
Being is the concept only as it is in itself. Its determinations have being, i.e.
in their difference they are others opposite one another, and their further
determination (the form of the dialectical) is a process of passing over into
an other. This progressive determination is at once a matter of settingforth
and thereby unfolding the concept, as it is in itself, and at the same time
the process of being entering into itself, a deepening of it within itself The
explication of the concept in the sphere of being becomes the totality of
being, precisely to the extent that the immediacy of being or the form of
being as such is sublated in the process.
§ 85
Being itself as well as the subsequent determinations, not only those of
being but also the logical determinations in general, can be regarded as
the definitions of the absolute, as metaphysical definitions of God. More
specifically, only the first simple determination within a given sphere, and
then the third, which is the return from a difference to the simple relation to
itself, can always be regarded in this way. For, to define God metaphysically
means to express his nature in thoughts as such. But logic comprises all
thoughts as they are while still in the form of thoughts. By contrast, the
second determinations, making up a given sphere in its difference (Diffirenz] ,
are the definitions of the finite. But if the form of definitions were used,
this would entail envisaging a representational substratum. For even the
absolute, what is supposed to express God in the sense and in the form
of thought, remains merely an intended thought, i.e. a substratum that as
such is indeterminate, relative to its predicate as the determinate and actual
expression in thought. Because the thought, the basic matter solely at issue
here, is contained only in the predicate, the form of a proposition, like that
subject, is something completely superfluous (c£ § 31 and the chapter on
the judgment below [§§ 166 et seq.]).
Addition. Each sphere of the logical idea proves to be a totality of determinations
and a presentation of the absolute, and so too does being, which includes within
itself the three stages of quality, quantity, and measure. Quality is, to begin with,
the determinacy that is identical with being in the sense that something ceases
to be what it is when it loses its quality. By contrast, quantity is. the determinacy
that is external to being and indifferent in relation to it. Thus, for instance, a
house remains what it is, whether it is bigger or smaller, and red remains red, be it
brighter or darker. The third stage of being, measure: is the unity of the first two,
qualitative quantity. All things have their measure: that is, they are quantitatively
determined, and their being either this big or bigger is indifferent to them. At
the same time, however, this indifference has its limits, and if those limits are
overstepped by an additional more or less, things cease to be what they were. From
measure there then results the progression to the second main sphere of the idea,
namely essence.
The three forms of being mentioned here are also the poorest, that is to say, the
most abstract, just because they are the first. The immediate sensory consciousness,
insofar as its behaviour involves thinking, is chiefly limited to the abstract deter-
minations of quality and quantity. This sensory consciousness is usually regarded
as the most concrete and thus also the richest. It is so, however, only in terms of
its material, whereas it is in fact the poorest and most abstract consciousness with
respect to the content of its thoughts.
A. QUALITY
a. Being
§ 86
Pure being constitutes the beginning, because it is pure thought as well as
the undetermined, simple immediate, and the first beginning cannot be
anything mediated and further determined.
All the doubts and reminders that might be raised against beginning
the science with abstract, empty being take care of themselves
through the simple consciousness of what is implied by the nature
of a beginning. Being can be determined as 'I = 1', as the absolute
indijfirence or identity, etc. In the need to begin with something
absolutely certain, i.e. the certainty of oneself, or with a definition or
intuition of the absolutely true, these and other similar forms can be
regarded as what must be the first. However, insofar as mediation is
,
137
already present within each of these forms, they are not truly the
first. Mediation means to have gone from a first to a second
and to emerge from something differentiated [Heroorgehen aus
Unterschiedenen]. If'I = I' or even the intellectual intuition is
genuinely taken as simply the first, then in this pure immediacy it
is nothing else but being, just as, conversely, pure being, insofar as it
is no longer this abstract being, but being that contains mediation
within itself, is pure thinking or intuiting.
When being is expressed as a predicate of the absolute, this
provides the first definition of the latter: the absolute is being. This
is (in the thought) the absolutely first, most abstract, and most
impoverished definition. It is the definition of the Eleatics, but at the
same time also the familiar one that God is the sum total [InbegriffJ of
all realities. The point is that one is supposed to abstract from the
limitedness inherent in every reality, so that God is nothing but the
real in all reality, the supremely real. Insofar as reality already contains
a reflection, this idea is expressed more immediately in what Jacobi
says about the God of Spinoza, namely that he is the principium of
being in all existence.
Addition I. When beginning with thinking, we have nothing but thought in
the sheer absence of any determination of it [in seiner reinen Bestimmungslosigkeitl,
since for a determination one and an other are required. In the beginning, however,
we have as yet no other. The indeterminate [Bestimmungslose], as we have it here,
is the immediate, not the mediated absence of determination, not the sublation of
all determinacy; but the immediacy of the absence of determination, the absence
of determination prior to all determinacy, the indeterminate as the very first. But
this is what we call 'being'. It is not to be sensed, intuited, or represented; instead
it is the pure thought, and as such it constitutes the beginning. Essence, too, is
something indeterminate, but the indeterminate that, having gone through the
mediation, contains within itself the determinacy as already sublated.
Addition 2. We find the various stages of the logical idea in the history of
philosophy, in the shape of philosophical systems that have successively emerged,
each of which has a particular definition of the absolute as its foundation. Now
just as the unfolding of the logical idea proves to be a progression from the
abstract to the concrete, so, too, the earliest systems in the history of philosophy
are the most abstract and thus at the same time also the most impoverished. The
relationship of the earlier to the later philosophical systems is, generally speaking,
the same as the relationship of the earlier to the later stages of the logical idea
and, to be sure, in such a way that the later ones contain within them the earlier
ones as sublated. This is the true meaning of the refutation of one philosophical
system by another, and more specifically of the earlier by the later system, a
common occurrence in the history of philosophy that is so otten misunderstood.The Encyclopedia Logic The Encyclopedia Logic
When the refutation of a philosophy is discussed, this tends at first to be taken
merely in an abstractly negative sense, such that the refuted philosophy has no
validity whatsoever anymore, that it has been discarded and done away with. If
this were so, the study of the history of philosophy would have to be regarded as an
altogether sad business, since study of it teaches how all philosophical systems that
have appeared over time have been refuted. However, just as one must admit that
all philosophies have been refuted, it must also be maintained that no philosophy
has ever been refuted or is even capable of being refuted. The latter is the case
in rwo connections, on the one hand, inasmuch as every philosophy worthy of
the name has the idea as such for its content, and on the other, inasmuch as·
each philosophical system has to be regarded as the presentation of a particular
moment, or a particular stage in the process of the development of the idea. Hence,
refuting a philosophy merely means that its limitation has been transcended and
its particular principle downgraded to an ideal moment. Accordingly, as far as its
essential content is concerned, the history of philosophy deals not with the past,
but with what is eternal and absolutely present, and its result must be compared
not to a gallery of errors of the human spirit, but rather to a pantheon of divine
figures [Gottergestalten l. These divine figures are the various stages of the idea as
they emerged successively in the dialectical development. Now it is left to the
history of philosophy to demonstrate in greater detail the extent to which the
unfolding of its contents that takes place in it agrees with the dialectical unfolding
of the pure, logical Idea, on the one hand, and diverges from it, on the other. All
that needs to be mentioned here is that the beginning of the logic is the same
as the beginning of the history of philosophy proper. We find this beginning in
the Eleatic philosophy, and more specifically in that of Parmenides who construes
the absolute as being when he says that 'only being is, and nothing is not'. This
is to be regarded as the proper beginning of philosophy because philosophy is,
generally speaking, a process of knowing by way of thinking [denkendes Erkennen],
but here for the first time pure thinking has been taken hold of and become an
object [gegenstandlichl for itself.
Human beings have thought from the beginning, to be sure, since they distin-
guish themselves from animals only through thinking. And yet it took thousands
of years before it came to grasping thought in its purity and at the same time as
absolutely objective. The Eleatics are famous for being bold thinkers. However,
this abstract admiration is often accompanied by the remark that these philoso-
phers nonetheless went too far by recognizing being alone as the true and denying
the truth of everything else that forms the object of our consciousness. Now it
is indeed perfectly correct to say that one must not stop at mere being. Still, it
is thoughtless to regard the remaining contents of our consciousness as existing
so to speak alongside and outside of being or as something that is there merely in
addition to it. By contrast, the true relationship here is that being as such is not
something fixed and ultimate but, rather, that it changes over dialectically into its
opposite, which, likewise taken immediately, is nothing. Thus it remains true in
the end that being is the first pure thought, and that whatever else may be made
the beginning (whether the 'I = 1', the absolute indifference, or God himself), it is at first only something represented and not something thought, and that in
terms of its thought contents it is only being after all.
139
§ 87
Now this pure being is a pure abstraction and thus the absolutely negative
which, when likewise taken immediately, is nothing.
The second definition of the absolure, namely that it is nothing,
followed from this. This conclusion is, indeed, entailed by saying
that the thing-in-itself is the undetermined, utterly devoid of form
and therefore of content. So, too, if it is said that God is simply the
supreme being and nothing else, then he is being declared, as
such, to be the very same negativity. The nothing that Buddhists
make the principle of everything and the ultimate end and goal of
everything is the same abstraction. - 2. When the opposition is
expressed in this immediate way as one of being and nothing, it
seems all too evident that it is null and void for one not to try to fix
[upon some determinate sense of] being and to save it from this
transition. In this respect, thinking the matter over is bound to fall
prey to looking for a fixed determination for being through
which it would be differentiated from nothing. For instance, one
may take it to be what persists in all change, the infinitely
determinable matter and so forth, or again, without thinking it
through, to be any given individual concrete existence [einzelne
Existenz]' the next best sensory or spiritual entity. However, none
of these further and more concrete determinations leave being as
pure being, as it is here immediately in the beginning. It is nothing
only in and because of this pure indeterminacy, something
inexpressible; its difference from nothing is a mere opinion [eine
blojJe Meinung]. - We are concerned here exclusively with the
consciousness of these beginnings, namely that they are nothing
but these empty abstractions and that each of them is as empty as
the other. The drive to find in being or in both a fixed meaning is
the very necessity that expands [weiterfohrt] being and nothing and
gives them a true, i.e. concrete meaning. This development is the
logical elaboration and the progression presented in what follows.
The process of thinking them over that finds deeper determinations
for them is the logical thinking by means of which these
determinations produce themselves, not in a contingent but in
I.The Encyclopedia Logic
The Encyclopedia Logic
a necessary manner. Each subsequent meaning they receive is
therefore to be regarded only as a more specific determination and a
truer definition of the absolute. Such a definition will then no longer
be an empty abstraction like being and nothing, but rather
something concrete in which both being and nothing are
moments. - The highest form of nothingness for itself would be
freedom, but freedom is the negativity that deepens itself within itself
to the point of the utmost intensity and is itself affirmation, and
absolute affirmation at that.
Addition. Being and nothing are at first only supposed to be distinguished,
i.e. their difference is at first only in itself, but not yet posited. If we talk about
a difference at all, then we have two and in each case a determination not to
be found in the one applies to the other. But being is absolutely devoid of all
determination, and nothing is the very same lack of determination. The difference
berween these rwo is therefore only intended the totally abstract difference that
is at the same time no difference. In all other cases of distinguishing we always also
have something common that subsumes the distinct items under it. For instance,
when we speak of rwo different genera, then the genus is what is common to
both. Similarly, we say there are natural and spiritual essences. Here, the essence
is something that belongs to both. In the case of being and nothing, however,
the difference is bottomless, and precisely for that reason there is none, for both
determinations represent the same bottomlessness. Suppose one wanted to say,
for instance, that both are after all thoughts, and hence thought is common to
bOth. One would then overlook the fact that being is not a specific, determinate
thought but rather the as yet entirely undetermined thought, and for that very
reason indistinguishable from nothing. - Again, being may also be represented as
absolutely rich and nothing as absolutely poor. But when we regard the entire world
and say of it that everything is and nothing further, we leave all determinateness
aside and instead of absolute fullness we only retain absolute emptiness. The same
comment can be made about its application to the definition of God as mere
being. Standing over and against this definition with equal justification is the
Buddhist definition that God is nothingness, with its implication that a human
being becomes God through self-annihilation.
§ 88
Conversely, nothing, as this immediate, self-same [category], is likewise the
same as being. The truth of being as well as of nothing is therefore the unity
of both; this unity is becoming.
1. The proposition 'Being and nothing are the same' appears to be 'such
a paradoxical proposition for the representation or the understanding
that one perhaps believes that it is not meant seriously. And indeed it
is one of the hardest thoughts that thinking imposes upon itself, for
2.
being and nothing are the opposite in its complete immediacy, that is to
say, without there already being posited in one of them a determination
that would contain its relation to the other. And yet, they do contain this
determination, as has been demonstrated in the previous section, namely,
the determination that is the same in both. The deduction of their unity
is thus entirely analytical, just as in general the whole progression in
philosophizing (insofar as it is a methodical, i.e. a necessary progression)
is nothing other than merely the positing of what is already contained
in a concept. But as correct as the unity of being and nothing is, so
it is also correct that they are absolutely different, i.e. that the one is not
what the other is. However, since at this point the difference has not yet
become determinate (for being and nothing are still what is immediate),
how it bears on them is something that cannot be said, it is something
merely meant [die blofe Meinung].
It does not require a great deal of wit to ridicule the proposition that
being and nothing are the same, or rather to bring up absurdities with
the false assurance that they are the consequences and applications of
it; for example, that according to that proposition it would be the same
whether my house, my assets, the air we breathe, this city, the sun,
right, spirit, God are Of not. For one thing, in examples such as these,
particular purposes or the utility something has for me ate surreptitiously
introduced, and it is asked whether it makes no difference to me, if the
useful thing exists or not. Philosophy is indeed just the doctrine that is
meant to liberate man from an infinite number of finite purposes and
goals, and to make him indifferent to them such that it is indeed all
the same to him whether such things are or not. But generally speaking,
as soon as we are talking about some contents, a connection is thereby
posited with other concretely existing things, purposes, etc. that ate pre-
supposed as valid, and it is then made dependent on such presuppositions,
whether the being or not-being of a determinate content is the same or
not. A difference foil of content is surreptitiously substituted for the
empty difference between being and nothing. - But for another thing,
purposes that are in themselves essential, absolute concrete existences
[absolute Existenzen] and ideas are placed under the determination of
being or not-being. Such concrete objects are something quite differ-
ent from mere beings or not-beings; poor abstractions such as being and
nothing (which are the poorest of all just because they are the determi-
nations only of the beginning) are completely inadequate to the nature
of those objects; a genuine content has long since transcended these
abstractions themselves and their opposition. - In general, if something
concrete is surreptitiously substituted for being and nothing, the usual
thing happens to this thoughtlessness, namely it entertains and talks
about something quite different from what is at issue. And what is at
issue here is merely abstract being and nothing.
3. It can easily be said that one does not comprehend the unity of being and
nothing. The concept of it, however, was stated in the preceding sections,
and it is nothing over and above what has been stated. Comprehending
it means nothing other than apprehending this. But by 'comprehend-
ing', something broader than the concept proper is understood. A more
manifold, richer consciousness, a representation is demanded, with the
result that a concept of this sort is put forward as a concrete case with
which thinking in its ordinary routine would be more familiar. To the
extent that the inability to comprehend expresses only that one is unac-
customed to holding onto abstract thoughts without any sensory input
and to grasping speculative sentences, there is nothing further to be said
than this, namely that philosophical knowledge [Wissen] is indeed of
a different sort from the kind of knowledge one is accustomed to in
ordinary life, as it also is from what reigns in other sciences. If, however,
the inability to comprehend means only that one is unable to repre-
sent this unity of being and nothing to oneself, then this is in fact so
little the case that to the contrary everybody possesses infinitely many
representations of this unity. That one does not possess such represen-
tations can mean only that one fails to recognize the concept under
discussion in any of those representations and that one does not know
that they are examples of it. The example that comes most readily to
mind is that of becoming. Everybody has a representation of becoming
and will equally admit that it is one representation; further, that when
one analyses it the determination of being, but also that of its absQlute
other, nothing, is contained therein; furthermore, that these two deter-
minations exist undivided in this one representation, so that becoming
is thereby the unity of being and nothing. - Another example equally
ready to hand is that of the beginning. The basic matter is not yet in
its beginning, but the beginning is not merely its nothing either; rather
being is already contained therein. The beginning is itself also a becom-
ing, but it already expresses the relation to the further progression. -
If one wanted to follow the usual procedure of the sciences, one might
start the Logic with the representation of the beginning thought in its
purity, i.e. with the beginning qua beginning, and to analyse this rep-
resentation. Perhaps one would then more easily accept as the result of
this analysis that being and nothing show themselves as undivided in a
single thought [in Einem ungetrennt]. 4. In addition, we must further note that the expressions 'Being and noth-
143
ing are the same' or 'the unity of being and nothing' and similarly all
other such unities (e.g. that of subject and object [Objekt], and so on)
are rightly objectionable. The awkwardness and incorrectness lies in
the fact that the unity is emphasized, and while the difference [Ver-
schiedenheit] is indeed contained in it (because the unity posited is one
of being and nothing, for instance), this difference is not simultaneously
stated and acknowledged. Instead, it seems that one is merely abstract-
ing illegitimately from it and not taking it into consideration. Indeed, a
speculative determination cannot properly be expressed in the form of
. such a proposition: unity is supposed to be articulated in the difference
that is simultaneously present and posited. As their unity, becoming is
the true expression of the result of being and nothing. It is not only
the unity of being and nothing, but the unrest in itself - the unity that
as relation to itself is not merely immobile, but is within itself against
itself on account of the difference of being and nothing contained in
it. - Existence [Dasein] is, by contrast, this unity, or becoming in this
form of unity; this is why existence is one-sided and finite. It is as if the
opposition had disappeared; It is contained in the unity only in itself,
but not posited in the unity.
5. Standing in contrast to the proposition that being is the transitioning
into nothing and nothing the transitioning into being (this being the
principle of becoming) is the proposition that 'Nothing comes from noth-
ing' or 'something can only come from something', i.e. the proposition
of the eternity of matter, pantheism. The ancients made the simple
reflection that the proposition 'something comes from something' or
'nothing comes from nothing' does indeed sublate becoming. For that
out of which something comes to be and that which comes to be are one
and the same. There is nothing here but a proposition of the identity
of the abstract understanding. It must strike one as curious, however, to
see the propositions 'nothing comes from nothing' or 'something comes
only from something' put forward quite na'ively even in our times with
neither any awareness that they are the foundation of pantheism, nor any
familiarity with the fact that the ancients considered these propositions
quite exhaustively.
Addition. Becoming is the first concrete thought and thus the first concept,
whereas being and nothing are empty abstractions. When we talk about the
concept of being, the latter can consist only in becoming, since as being it is the
empty nothing and as such the empty being. In being, then, we have nothing and
in it being. This being, however, that persists in being with itself in nothing isThe Encyclopedia Logic The Encyclopedia Logic
becoming. In the unity of becoming, the difference [Unterschiedj must not be left
out, for without it one would return to abstract being. Becoming is merely the
positedness [d:ts Gesetztseinl of what being truly is.
One very often hears the claim that thinking is opposed to being. In the face of
such an affirmation, however, it should first be asked what we are to understand
by being. When we take up being as it is determined by reflection, the only thing
we can say about it is that it is the absolutely identical and affirmative. If we then
consider thinking, it cannot escape us that at the very least it is likewise absolutely
identical to itsel£ To both being and thinking, then, the same determination
applies. This identity of being and thinking must, however, not be taken in a
concrete sense, and hence one is not to say that a stone that has being is the
same as a thinking human being. Something concrete is quite different from the
abstract determination as such. But in the case of being, there is no talk of anything
concrete, for being is precisely what is entirely abstract. Accordingly, the question
concerning the being of God who is in himself infinitely concrete, is also of little
interest.
As the first concrete determination of thought, becoming is also at the same time
the first true determination of thought. In the history of philosophy, it is the system
of Heraclitus that corresponds to this stage of the logical Idea. When Heraclitus
says 'Everything is in flux' (1Tano pEl), becoming is thereby pronounced to be
the fundamental determination of all there is, whereas the Eleatics by contrast, as
mentioned earlier, construed being alone rigid being, devoid of any process - as
true. With teference to the principle of the Eleatics Democritus later comments:
'Being is no more than not-being' (OOOEV \..Iai\i\ov ,0 OV ,OV \..IT] QV,Os seri).
He thereby expresses the negativity of abstract being and its identity, posited in
becoming, with a nothing that is equally untenable in its abstraction. - At the
same time we have here an example of the true refutation of one philosophical
system by another, a refutation that consists precisely in exhibiting the dialectic
of the principle of the refuted philosophy and in downgrading this principle to
an ideal moment in a higher, more concrete form of the idea. - But furthermore,
becoming, too, is in and for itself as yet a supremely impoverished determination
that has to funher deepen and fulfil itself in itself. We have such a deepening
of becoming within itself in, for instance, life. The latter is a becoming, but its
concept is not exhausted by this. We find becoming in an even higher form in
spirit. Spirit is likewise a becoming, but a more intensive, richer one than the
merely logical becoming. The moments whose unity is spirit are not the mere
abstractions of being and nothing, but the system of the logical idea and nature. contradiction, becoming collapses into the unity in which both are sub-
lated. Its result is therefore existence.
I44
b. Existence
§ 89
The being in becoming, as one with nothing, and the nothing that is
likewise one with being are only vanishing [moments]. Due to its inner
I45
In connection with this initial example, we are once and for all
to be reminded of what was stated in § 82 and in the Remark. What
alone can ground [begriinden] a progression and a development in
knowing [Wissen] is to hold on to the results in their truth. Suppose
a contradiction is pointed up in any sort of object or concept (and
there is simply nothing anywhere in which a contradiction, Le.
opposite determinations, could not and would not have to be
pointed out, for the understanding's process of abstracting vi?le!ltly
holds on to one determinacy, while striving to obscure and ehmmate
the consciousness of the other determinacy that is contained in it).
When such a contradiction is recognized, the conclusion is usually
drawn that' Therefore, the object is nothing, just as Zeno first
demonstrated with regard to movement, namely that it contradicts
itself and that therefore it does not exist, or as the ancients
recognized the two kinds of becoming, namely coming-to-be and
passing away, to be untrue determinations by stating that the C!ne,
i.e. the absolute, neither comes into being nor passes away. ThlS
kind of dialectic thus merely stops at the negative side of the result
and abstracts from what is at the same time actually on hand,
namely a determinate result, here a pure nothing, but a nothing that
contains being and likewise a being that contains nothing within
itself. Thus, existence is (I) the unity of being and nothing in
which the immediacy of these determinations has disappeared and
with it the contradiction in their relationship, - a unity in which
they are now only moments. (2) Since the result is the sublated
contradiction, it is in the form of a simple unity with itself or itself as
being, but a being with negation or determinateness. It is becoming
posited in the form of one of its moments, that of being.
Addition. It is also contained in our representation that when there is becom-
ing, something comes out of it and that therefore becoming h~ a result. But
there then arises the question how becoming manages not to remam mere becom-
ing but to have a result. The answer to this qu~stion de:ives ~ro~ hat ab~ve
has shown itself to us as becoming. For becommg contams Within Itself being
and nothing, and in such a way that these twO change over into one another
absolutely and mutually sublate each other. In this way, becoming .prove~ itse!f
to be what is restless through and through, yet unable to preserve Itself m thiS
abstract restlessness. For because being and nothing disappear in becoming and
its concept consists in this alone, becoming is thus itself something vanishing,
:vThe Encyclopedia Logic The Encyclopedia Logic
like a fire that extinguishes itself by consuming its materiaL The result of this
process, however, is not empty nothing but being that is identical with negation
something we call existence and whose meaning for now proves to be this: to have
become. being-for-another- a breadth of existen~, of something. The being of qual-
ity as such, as opposed to this relation to something other, is being-in-itself.
Addition. The foundation for every determinacy is negation (omnis determinatio
est negatio [all determination is negation], as Spinoza says). Thoughtless opining
regards determinate individual things alone as positive and fastens on to them
under the form of being. Nothing much is accomplished, however, with mere
being, for as we saw earlier it is what is absolutely empty and insubstantial.
Incidentally, this much is correct concerning the confusion mentioned here of
existence as being that is determined [Dasein als bestimmtes Sein] with abstract
being, namely that the moment of negation is, indeed, still contained merely in
a veiled state, as it were, in existence, while the moment of negation emerges
freely only in being~for~itself and there assumes its rightful position. - If now
we consider existence also as a determinacy that simply is, we then have what is
understood by reality. In this way one ralks of the reality of a plan, for instance,
or of an intention, and understands by it that such things are no longer merely
something inner, subjective, but instead have emerged into existence. In the same
sense, the body may then be called the reality of the soul, and this particular right
the reality of freedom or, quite generally, the world may be called the reality of the
divine concept. In addition, however, there is also talk of reality in still another
sense, where what is understood by it is that something behaves in accordance
with its essential determination or its concept. This happens, for instance, when
it is said 'this is a real [reei~ occupation' or 'this is a real [ree/i] human being'. In
these cases it is a matter not of the immediate, external existence, but instead of
the agreement of an existent [eines Daseienden] with its concept. So construed,
however, reality is not that different from ideality, which we will initially come to
know as being-for~itself.
§ 90
(0) Existence [Dasein] is being with a detenninacy that is immediate or
that simply is, i.e. quality. Existence qua reflected into itself in this its
determinacy is an existent [Daseiendes], something [Etwas]. - The categories
that develop in connection with existence need to be specified in a summary
fashion only.
Addition. Generally speaking, quality is the determinacy that is identical with
being and immediate, in contrast with quantity that is to be considered next.
Quantity is, to be sure, likewise a determinacy of being, but one that is no longer
identical with it. Quantity is instead a determination indifferent to being and
external to it. - Something is what it is by virtue of its quality, and when it loses
its quality it stops being what it is. Moreover, quality is essentially a category
merely of the finite. For this reason, it has its proper place only in nature, not
in the spiritual world. Thus, for instance, in nature the so-called simple types of
matter, e.g. oxygen, nitrogen, etc., are to be considered concretely existing qualities
[existierende Qualitiiten]. By contrast, in the sphere of spirit quality occurs only
in a subordinate manner and not in such a way that any given determinate shape
of spirit would be exhaustively characterized by means of it. For instance, when
we consider subjective spirit, the object of psychology, it is indeed possible to say
that the logical meaning of what one calls character is that of quality. But this
is not to be understood as though character were a determinacy that penetrates
the soul and is immediately identical with it as is the case with the simple types
of matter in nature mentioned above. Quality, however, shows itself in a more
determinate manner even in connection with spirit to the extent that the latter is
in an unfree, sick condition. This is notoriously the case with the state of passion
and with passion that has escalated to madness. It can fittingly be said of a mad
person whose consciousness is completely pervaded by jealousy, fear, etc., that his
consciousness is determined as a quality.
§ 91
As a determinacy that simply is [seiende Bestimmtheit] over against the nega-
tion that is contained in it but distinct from it, quality is reality [Realitiit].
Negation, no longer as the abstract nothing but as an existent and some-
thing, is only the form in the latter, it is as being-other. Because this being-
other is its own determination, but at first distinct from it, quality is
147
§92
(13) The being that is fastened onto as distinct from determinacy, i.e. the
being~in-itself, would be merely the empty abstraction of being. In existence,
determinacy is one with being, and at the same time posited as negation,
i.e. limit, barrier. Being other is thus not something indifferent outside of
it but instead its own moment. By virtue of its quality, something is, first,
finite and, second, alterable, so that finitude and alterability belong to its
being.
Addition. In existence, negation is still immediately identical with being, and it
is negation that we call a limit. Something is what it is only within its limit and
due to its limit. Hence one must not regard the limit as something that is merely
external to existence; rather it permeates existence as a whole. The construal of the
limit as a merely external determination of existence is due to the conflation of
the quantitative with the qualitative limit. At issue here is for now the qualitative
limit. If we consider, for instance, a plot of land that is three acres, this is thenThe Encyclopedia Logic The Encyclopedia Logic
its quantitative limit. In addition, however, this plot of land is also a meadow
and not a forest or a pond, and this is its qualitative limit. - Insofar as human
beings want to be actual, they must exist [mul dasein] and to this end they must
limit themselves. Those who are too dismayed at the finite do not accomplish
anything actual, but instead remain trapped in the abstract and fade away into
themselves.
When we now consider more closely what we have here in the case of the
limit, we find that it contains in itself a contradiction and thus proves itself to
be dialectical. For, on the one hand, the limit constitutes the reality of existence,
but on the other hand it is the negation of the latter. Moreover, however, as the
negation of [the] something the limit is not an altogether abstract nothing, but a
nothing that is [ein seiendes Nichts] or what we call an other. When thinking of
the something, the [concept of the] other immediately comes to mind, and we
know that there is not only something but also an other as well. But the other
is not just something that we simply find such that the something could also
be thought without it. Rather, something is in itself the other of itself, and in
the other the limit of the something becomes objective for ii:. When we now ask
about the distinction between the something and the other, it is evident that both
are the same, an identity that is expressed in Latin by the designation of both
as aliud - aliud. The other opposed to the something is itself a something, and
accordingly we say 'something else' [etwas Anderes: lit. 'something other']. So, too,
the first something is in turn itself an other vis-a-vis the other that is likewise
determined as a something. When we say 'something else', we at first imagine that
the something, taken by itself, is only something and the determination of being
an other accrues to it on account of an external consideration alone. Thus, for
instance, we think that the moon, which is something other than the sun, could
also exist even if the sun did not. In fact, however, the moon (as a something) has
its other in and of itself [an ihm selbst] and this constitutes its finitude. Plato says:
'God made the world from the nature of the One and the Other (TOU ETepov);
these he brought together and out of them fashioned a third which is of the nature
of the One and the Other'.20 - With this, the nature of the finite is being expressed
as such, which qua something does not stand over against the other indifferently,
but is in itself the other of itself and in this way alters itself. In the alteration
the inner contradiction shows itself with which existence is intrinsically beset and
which drives it beyond itself. Existence at first appears to the representation as
simply positive and at the same time as remaining tranquilly within its boundary.
To be sure, we also know that all finite things (and such is existence) are subject to
alteration, but this alterability of existence appears to the representation as a mere
possibility, the realization of which is not grounded in itself. In fact, however,
it is part of the concept of existence to alter itself, and alteration is merely the
manifestation of what existence is in itself. Living things die, and they do so
simply because they carry the germ of death in themselves. § 9J
2?
Translators' note: see the two elements of the indivisible and the divisible in Timaeus 34-5.
149
Something becomes an other, but the other is itself a something, hence it
likewise becomes an other, and so on and so forth ad infinitum.
§ 94
This infinity is the bad or negative infinity in that it is nothing but the
negation of the finite, which, however, re-emerges afresh and thus is just
as much not sublated. In other words, this infinity expresses only that
the finite ought to be sublated. The progression to infinity stops short at
expressing the contradiction that is contained in the finite, namely that
it is something as well as its other and that it is the perpetual continuance
of the alternation of these determinations each of which brings about the
other.
Addition. When we let the moments of existence, namely something and the
other, fall apart, we have the following: something becomes an other, and this
other is itself a something that then as such likewise alters itself, and so on ad
infinitum. Reflection believes it has reached something very lofty here, indeed
even the loftiest [thought]. This progression to the infinite is, however, not the
true infinite. The latter consists, rather, in being with itself in its other, or, put in
terms of a process, to come to itself in its other. It is of great importance to grasp
the concept of the true infinity properly and not merely to stop short at the bad
infinity of the infinite progression. When the infinity of space and time is under
discussion, it is at first the infinite progression that one tends to focus on. Thus
one says, for instance, 'this time', 'now', and this boundary is then continuously
surpassed, backwards and fotwards. It is the same with space about whose infinity
edifying astronomers put forth many empty declamations. It is the~ ~so :rricall.y
asserted that thinking must give up when it starts to contemplate thiS In.fifllty. ThiS
much is indeed correct, namely that we eventually abandon proceedIng further
and further in such contemplation, but on account of the tediousness, not the
sublimity, of the task. Engaging in the contemplation of this infinite progression
, is tedious because the same thing is incessantly repeated here. A limit is posited,
it is surpassed, then again a limit, and so on endlessly. So there is nothing here
but a superficial alternation that remains stuck in the finite. If it is thought that
through stepping forth into that infinity one liberates oneself from the finite: then
this is indeed merely the liberation of fleeing. The one who flees, however, IS not
yet free, for in fleeing he is still dependent on what he flees. If it is then further said
that the infinite cannot be reached, then this is quite right, but only because the
determination of being something abstractly negative is read into it. Philosophy
does not waste its time with such empty and merely transcendent Uemeitigen]
things. What philosophy deals with is always something concrete and ~bsolutel?,
present. - The task of philosophy has occasionally been framed by sayIng that ItThe Encyclopedia Logic
must answer the question of how the infinite resolves to move beyond itself. To
thisquestion, which is predicated on the fixed opposition between the infinite and
the finite, one can only answer that this opposition is something untrue and that
the infinite is indeed eternally beyond itself and also eternally not beyond itself.
Moreover, when we say the infinite is the not-finite, we have thereby indeed
already uttered the truth, for since the finite is the first negative, the not-finite is
the negative of negation, Le. the negation that is identical with itself and thus at
the same time true affirmation.
The infinity of reflection here under discussion is merely the attempt to reach the
true infinity; [in other words, it is] a hapless hybrid. Generally speaking, this is the
standpoint of philosophy that has been maintained and upheld [geltend gemacht]
in Germany in recent times. The finite here merely ought to be sublated, and the
infinite ought to be not merely something negative, but something positive as well.
This ought always carries within itself the impotence of recognizing something
as legitimate that nonetheless cannot maintain and uphold itself. With respect
to ethics, the Kantian and the Fichtean philosophy have stopped short at this
standpoint of the ought. The perennial approximation to the law of reason is the
utmost that can be achieved on this path. The immortality of the soul was then
also based on this postulate.
§ 95
(y) What is in fact the case is that something becomes an other and
the other generally becomes something other. In the relation to an other,
something is itself already an other opposite it. Hence, since that into
which it makes the transition is entirely the same as that which makes the
transition (both have no further determination than this, which is one and
the same, the determination to be an other), something comes together
only with itselfin its transition into something other, and this relation to
itself in its transition and in the other is the trne infinity. Or, considered
negatively, what is altered [veriindert] is the other [das Andere]; it becomes
the other of the other. In this way, being but as negation of negation - is
re-established and is being-for-itself
The dualism that makes the opposition of the finite and the infinite
insuperable fails to make the simple observation that in this manner
the infinite is at once only one of the two, that it is thus made into
merely one particular for which the finite is the other particular.
Such an infinite that is only a particular, next to the finite which
makes up its boundary and limit, is notwhatit is supposed to be;
not the infinite, but merely finite. - In such a relationship, where the
finite is hither, the infinite thither, the one placed on this side, the
The Encyclopedia Logic
other on the other side, the finite is accorded the same honour of
subsisting and being self-standin'g [Bestehen und Selbstiindigkeit] that
the infinite is. The being of the finite is made into an absolute being.
In such a dualism, it stands firmly for itself. If it were touched by the
infinite, so to speak, it would be annihilated. But it is supposed to be
untouchable by the infinite. There is supposedly an abyss, an
insurmountable chasm between the two, with the infinite remaining
absolutely on the other side and the finite on this side. While one
may believe that the assertion that the finite persists steadfastly
opposite the infinite gets one beyond all metaphysics, it in fact
stands squarely on the grounds of the most ordinary metaphysics of
the understanding. The same thing happens here which is expressed
by the infinite progression. At one moment, it is admitted that the
finite does not exist in and for itself that it is not a self-standing
actuality, not an absolute being, that it is only transitoty. The next
moment, this is immediately forgotten and the finite is represented
as existing entirely over against the infinite, absolutely separated from
it and exempted from annihilation, as self-standing and persisting
for itselE - While such thinking believes that it is elevating itself to
the infinite in this manner, the opposite happens to it - it arrives
at an infinite that is merely finite, and, instead ofleaving the
finite behind, permanently holds onto it, making it into an
absolute.
Based on these considerations concerning the emptiness
[Nichtigkeit] of the understanding's opposition of the finite and the
infinite (one may benefit from comparing Plato's Philebus with it),
it is easy to lapse into the expression that therefore the infinite and
the finite are one, that the true, i.e. true infinity, is determined and
declared to be the unity of the infinite and the finite. It is true that
phrasing the matter in such a way is in some sense correct, but it is
equally skewed and false (as was mentioned earlier with regard to
the unity of being and nothing). Furthermore, it invites the just
reproach of having finitized the infinite, the reproach of a finite
infinite. For in the above phrasing the finite appears as if untouched,
i.e. it is not explicitly stated that the finite has been sublated. - Or,
when one reflects that the finite, in being posited as one with the
infinite, could not indeed remain what it was outside this unity, and
that it would suffer at least some modification in its determination
(just as an alkali combined with acid loses some of its properties),
then the same thing should happen to the infinite, which, as theThe Encyclopedia Logic
negative, should have to be blunted by the other equally in turn.
And this is indeed what happens to the abstract, one-sided infinite
of the understanding. However, the true infinite does not behave
merely like the one-sided acid, but instead preserves itself Negation
of negation is not a neutralization. The infinite is the affirmative,
and only the finite is what is sublated.
In being-for-itself, the determination of ideality has made its
entry. Existence, construed at first only in terms of its being or its
affirmative nature, has reality (§ 91). Thus, too, finitude is at first
determined in terms of reality. But the truth of the finite is rather its
ideality. Likewise, the infinite of the understanding, which posited
next to the finite is itself merely one of the two finites, is something
untrue, something ideal [ideelles]. This ideality of the finite is the
chief proposition of philosophy, and every true philosophy is for
that reason idealism. The only thing that matters is not to take as the
infinite what is at once made into something particular and finite
in the determination of it. - This is why we have drawn attention
to this distinction here at some length. The fundamental concept
of philosophy, the true infinite, depends on this. This distinction
is taken care of by the very simple, and therefore perhaps
unremarkable, but irrefutable reflections contained in this
section.
c. Being-for-itself
§96
(0) Being-for-itself as relation to itself is immediacy, and as the relation
of the negative to itself it is a being that is for itself [Fursichseiendes], the
One what is in itself devoid of any distinction, hence, what excludes the
other from itself [das Andere aus sich Ausschliejfende].
:1ddition .. B~ing-for-itself is perfected quality and as such contains being and
ex.lSt~nce as.ltS Ideal moment~ within. it~el£ Qua being, being-for-itselfis the simple
relatIOn to Itself, and qua 'eXlstence It IS determined. This determinacy, however,
is no lo~ger the finite determinacy of something in its difference from the other,
but the Infinite determinacy that contains in itself the difference as sublated.
We ~a~e the most obvious example of being-for-itself in the I. To begin with,
qua eXisting we know ourselves to be distinct from other existents and related to
them. Furthermore, we know this expanse of existence to be at the same time
The Encyclopedia Logic
I53
sharpened, so to speak, into the simple form of being-for-itself. Saying T is the
expression of an infinite and at the same time a negative relation to oneself. It
can be said that human beings distinguish themselves from animals and hence
, from nature generally by knowing themselves [in each case] as an l. At the same
time, by this means, one expresses that natural things do not attain [the status
of] free being-for-itself Instead, by being confined to existence they are forever
merely being-for-another. - In addition, being-for-itself must be construed as
ideality generally, whereas existence, by contrast, was previously designated as
reality. Reality and ideality are often regarded as a pair of determinations standing
over and against one another, each with the same self-standing character, and it is
accordingly said that apart from reality there is also an ideality. However. ideality is
not something that there is apart from and alongside reality. Rather, the concept of
ideality consists expressly in being the truth of reality; that is to say, reality posited
as what it is in itself proves to be ideality. Hence, one must not believe that one
has accorded ideality the proper honour if one merely concedes that reality alone
does not suffice and that one must also acknowledge an ideality apart from reality.
An ideality such as this, along with or even above reality, would indeed be only an
empty name. Ideality has content only by being the ideality of something. This
something, however, is not merely an indeterminate this or that, but an existence
that is determined as reality and which possesses no truth, if taken in isolation.
It is not without reason that the difference between nature and spirit has been
construed in such a way that the former should be traced back to reality and
the latter to ideality as their fundamental determinations. Now nature is indeed
not something fixed and finished for itself, something that could therefore subsist
without spirit. Rather, nature achieves its end and truth only in spirit, and spirit
for its part is similarly not JUSt an abstract beyond of nature; rather, it exists and
validates itself as spirit only insofar as it contains in itself nature as sublated. We
are to be reminded here of the dual meaning of our German expression 'aufheben'
[to sublate]. By 'aufheben' we understand on the one hand something like dearing
OUt of the way or negating, and we accordingly speak of a law, for instance, or
an institution as having been 'aufgehoben'. On the other hand, however, aufheben
also means something like preserving, and in this sense we say that something is
well taken care of [gut aufgehoben, taken out of harm's way and put in a safe place].
This dual sense in linguistic usage according to which one and the same word has
a negative as well as a positive meaning must not be regarded as a coincidence or
even made the object of reproach to the language as causing confusion. Rather, in
it we should recognize the speculative spirit of our language that transcends the
either/or of mere understanding.
§ 97
(~) The relationship of the negative to itself is a negative relationship, hence
the distinguishing of the One from itself, the repulsion of the One, i.e. aThe Encyclopedia Logic
positing of many Ones. In accordance with the immediacy of that which
is a being-for-itself, these many are beings, and the repulsion of the Ories
that have being becomes in this respect their repulsion against each other
insofar as they are on hand, or a mutual excluding.
Addition. When we talk about the One, the first thing that tends to occur to us is
the Many. The question then arises where the Many come from. In representational
thought no answer is to be found to this question, since it considers the Many to be
immediately on hand, and since the One counts simply as one among the Many.
In terms of the concept, however, the One constitutes the presupposition for the
Many, and it is inherent in the thought of the One to posit itself as the Many. For
unlike being, the One as being-for-itself as such does not lack relatedness; rather,
it is a relation just as existence is. However, it does not relate as something does to
an other but instead, as the unity of something and an other, it is relation-to-itself,
and, furthermore, this relation is negative relation. With this, the One proves to
be what is absolutely incompatible with itself, what repels itself from itself, and
what it posits itself as is the Many. We may designate this side in the process of
being-for-itself with the figurative expression repulsion. One speaks of repulsion
first and foremost in considering matter, and one understands by it that, as a
Many in each one of these many Ones, matter behaves by excluding all the others.
Moreover, the process of repulsion must not be construed in such a way that the
One does the repelling and the Many are what is repelled. Rather, as was mentioned
earlier, the One is precisely just this, namely to exclude itself from itself and to
posit itself as the Many. Each of the Many, however, is itself a One and, because it
behaves as such, this ubiquitous repulsion changes over into its opposite, namely
attraction.
§ 98
(y) Of the Many, however, one is what the others are; each is a One as
well as one of the Many. They are therefore one and the same. Or, con-
sidered in itself, repulsion as the negative behaviour of the many Ones
to each other is equally essentially their relation to each other. And since
those to which the One relates in its repelling are Ones, it relates to
itself in them. Thus repulsion is equally essentially attraction, and the
excluding One or being-for-itself sublates itself. The qualitative deter-
minacy that has reached in the One its determinacy in-and-for-itself
has thus passed over into determinacy qua sublated, i.e. into being as
quantity.
The atomistic philosophy is the standpoint on which the absolute
determines itself as being-for-itself, as One, and as many Ones.
Repulsion, which shows itself in the concept of the One, has also
The Encyclopedia Logic
155
been assumed to be its fundamental force. Not, however, attraction
but coincidence, i.e. something thoughtless, is supposed to bring
them together. If one is fixated on the One as One, its coming
together with others must indeed be regarded as something quite
extrinsic. - The void that is adopted as the other principle in
addition to the atoms is repulsion itself, represented as the existing
nothing in between the atoms. The more recent atomism (and
physics continues to hold on to this principle) has given up atoms
insofar as it focuses on small particles, the molecules. In this, it has
drawn closer to sensory representation and abandoned thoughtful
determination. Moreover, insofar as a force of attraction is set
alongside the force of repulsion, the opposition has, it is true, been
made complete, and the discovery of this so-called force of nature has
been touted a lot. But the relationship of both to one another that
constitutes what is concrete and true about them would need to be
rescued from the state of cloudy confusion in which it has been left
even in Kant's Metaphysical Foundations o/the Natural Sciences.
In recent times, the atomistic approach has become even more
important in the political than in the physical sphere. According to
this view, the will of the individual as such is the principle of the
state. The attractive force is the particularity of the needs and
inclinations, and the universal, the state itself, is [based on] the
external relationship of the contr~ct.
Addition I. The atomistic philosophy constitutes an essential stage in the his-
torical development of the idea, and the principle of this philosophy generally is
being-for-itself in the shape of the Many. If today the atomistic doctrine is held in
high esteem even by those physicists who shun metaphysics, one should remember
here that one does not escape metaphysics, and more specifically the reduction of
nature to thoughts, by throwing oneself into the arms of atomism. For the atom
is indeed itself a thought, and hence the interpretation of matter as consisting of
atoms is a metaphysical interpretation. It is true that Newton explicitly warned
physics to guard against metaphysics. But to his credit it must be said that he
did not himself act by any means in accordance with this warning. Indeed, only
the animals are pure, unadulterated physicists, since they do not think, whereas
a human being as a thinking being is a born metaphysician. The only thing that
matters, therefore, is whether the metaphysics one applies is of the right kind,
namely whether, instead of the concrete logical idea, it is one-sided thought deter-
minations fixed by the understanding that one holds on to and that form the
basis of our theoretical as well as practical activities. This is the objection that
applies to the philosophy of atomism. As is often the case even today, the ancient
atomists regarded everything as a Many, and coincidence was then supposed toThe Encyclopedia Logic
The Encyclopedia Logic
bring together the atoms that floated around in the void. But the relationship of
the Many to each other is by no means merely accidental; this relationship has
its ground instead, as previously mentioned, in the Many themselves. It is Kant
who deserves the credit for having brought the way matter is construed to comple-
tion, by regarding it as the unity of repulsion and attraction. What this view gets
right is the fact that attraction must indeed be recognized as the other moment
contained in the concept of being-for-itself and that, as a consequence, attraction
belongs to matter just as essentially as repulsion does. But this so-called dynamic
construction of matter suffers from the defect that repulsion and attraction are
postulated as being on hand without further ado and are not deduced. From this
deduction the how and why of their merely alleged unity would have followed. By
the way, Kant explicitly insisted that one must regard matter not as on hand for
itself and equipped in passing, so to speak, with the two forces mentioned here,
but as obtaining instead only in their unity, and for a time German physicists went
along with this pure dynamics. In more recent times, and against the warning
of their colleague, the late Kastner, the majority of these physicists has found it
more comfortable to return to the atomistic standpoint and to regard matter as
consisting of infinitely small things called 'atoms'. These atoms are then supposed
to be set in relation to each other due to the play of the attractive, repulsive, and
whatever other forces that attach to them. This is then likewise a metaphysics and
one has, to be sure, quite sufficient reason to guard against it, given the lack of
thought in it.
Addition 2. The transition from quality to quantity indicated in the preced-
ing section is not to be found in our ordinary consciousness. The latter takes
quality and quantity to be a pair of self-standing determinations existing side
by side and it is accordingly said that things are not only qualitatively but also
quantitatively determined. Where these determinations come from and how they
relate to each other, these questions are not raised here. But quantity is noth-
ing other than quality sublated, and it is the dialectic of quality studied here by
virtue of which this sublation comes to pass. At first, we had being, and becoming
resulted as its truth. This formed the transition to existence whose truth we-saw
to be alteration. Alteration, in turn, showed itself in its result to be being-for-
itself that was exempt from the relation to an other and from its transition into
it. And, finally, being-for-itself proved to be the sublating of itself, and thus of
quality in general, in the totality of its moments on both sides of its process.
Now this sublated quality is neither an abstract nothing nor the equally abstract
and indeterminate being, but rather being that is indifferent to determinacy. It
is this shape of being that also surfaces in our ordinary representation as quan-
tity. Accordingly, we consider things first from the viewpoint of their quality,
and the latter we take to be the determinacy that is identical with the being of
the thing. As we proceed next to considering quantity, it offers us at once the
representation of an indifferent, external determinacy in the sense that, even if a
thing's quantity changes and it becomes greater or smaller, it still remains what
it is.
I57
B. QUANTITY
a. Pure quantity
§ 99
Quantity is pure being in which determinacy is posited as no longer one
with being itself, but as sublated or indifferent.
The expression magnitude is unsuitable for quantity, insofar as it
signifies first and foremost determinate quantity. 2. Mathematics
usually defines magnitude as what can be increased or decreased. As
faulty as this definition is (inasmuch as it repeats what is to be
defined [DeJinitum]), it still conveys this much, namely, that the
determination of magnitude is such that it is posited as alterable and
indifferent. Hence, apart from any alteration of it, e.g. an increase in
extension or intensity, the basic matter, for instance, a house or red,
does not cease to be a house or red. 3. The absolute is pure quantity.
This standpoint generally coincides with determining the absolute as
matter in which the form is indeed on hand, but as an indifferent
determination. Quantity also constitutes the basic determination
of the absolute, when it is grasped in such a way that in it (as the
absolutely indifferent) every distinction is only quantitative. - Pure
space, time, etc. may equally be taken as examples of quantity,
insofar as one is supposed to construe the real as an indifferent filler
of space or time.
1.
Addition. At first glance mathematics' customary definition of magnitude as
what can be increased or decreased seems to be more illuminating and plausible
than the conceptual determination contained in the above section. Looked at more
closely, however, it contains in the form of a presupposition and representation the
same [determination] as the concept of quantity that was the result of the logical
development. For, if it is said of magnitude that its concept consists in being able
to be increased or decreased, then it is stated precisely with this that magnitude
(or, more correctly, quantity, as distinct from quality) is a determination of the SOrt
that the specific basic matter behaves indifferently towards its alteration. As for
the earlier criticized deficiency in the customary definition of quantity, it consists
more specifically in the notion that increasing and decreasing mean nothing other
than determining a magnitude in different ways_ But if this were the case, quantity
would then be merely something alterable in general. But quality is alterable, too,
and the previously mentioned difference between quantity and quality is thenThe Encyclopedia Logic The Encyclopedia Logic
expressed as a matter of increasing or decreasing. This implies that the basic mat-
ter remains what it is, regardless of the direction in which the determination of
magnitude is changed. - At this point, it should also be noted that in philosophy
we are not at all concerned merely with correct definitions, much less with merely
plausible definitions, i.e. definitions whose correctness is immediately obvious to
representational consciousness. Rather, we are concerned with definitions that
have a proven record, i.e. definitions whose content has not merely been taken up
as something found, but one that is known to be grounded in free thinking and
thus at the same time known to be grounded in itself. This finds its application
in the current case in such a way that, however correct and immediately obvious
mathematics' customary definition of quantity might be, this would still not satisfy
the requirement of knowing to what extent this particular thought is grounded in
universal thinking and therefore necessary. There is a further consideration that is
linked to this point here. When quantity is taken up directly from representation
without being mediated by thinking, it easily happens that quantity is overesti-
mated with respect to its scope and even raised to an absolute category. This is
indeed the case when only those sciences whose objects can be submitted to a
mathematical calculus are recognized as exact sciences. Here that bad metaphysics
mentioned earlier (§ 98 Addition) shows itself again, replacing the concrete idea
with one-sided and abstract determinations of the understanding. Our knowing
would indeed be in bad shape, if, renouncing exact knowledge, we generally had
to be satisfied merely with a vague representation of such objects as freedom,
law, the ethical life, even God himself, merely because they cannot be measured
and calculated or expressed in a mathematical formula; and if, when it comes to
the more specific or particular derails of those matters, it would then be left to
each individual's whim to make of it what they want. - It is immediately obvious
what kind of pernicious practical consequences result from such a view. Looked at
more closely, the exclusively mathematical standpoint mentioned here (for which
quantity, this specific stage of the logical idea, becomes identical with the logical
idea itself) is none other than materialism. Indeed, this is fully confirmed in the
history of scientific consciousness, notably in France since the middle of the last
century. The abstractness of matter is precisely this: that the form is indeed on
hand in it, but merely as an indifferent and external determination. Inciden-
tally, the remarks added here would be greatly misunderstood, if one intended to
construe them as detracting in any way from the dignity of mathematics, or as if
by designating quantitative determination as merely external and indifferent, they
were supposed to encourage lethargy and superficiality, as though one could set the
quantitative determinations aside or at least that it was thus not necessary to rake
them seriously. Quantity is in any case a stage of the idea to which justice must be
done, initially as a logical category, but then also in the objective fgegenstand/ich]
world, the natural as well as the spiritual world. But here the difference between
them also becomes at once apparent, namely that the determination of magnitude
,is not of equal importance with respect to the objects of the natural and the spir-
itual world. For in nature, taken as the idea in the form of otherness and at the
same time of being-outs ide-itself, quantity is - precisely for that reason - of greater importance than in the world of spirit, which is the world of free interiority. To be
sure, we consider the spiritual content f{om the quantitative viewpoint as well, but
it is immediately obvious that when we contemplate God as a trinity the number
three has a much more subordinate significance than if we were to contemplate
the three dimensions of space, not to mention the three sides of a triangle whose
basic determination is just this, namely to be a surface delimited by three sides.
Furthermore, the difference mentioned between a greater and lesser imporrance
of the quantitative determination i~ also found in nature and, indeed, in the sense
that quantity plays, so to speak, a more important role in inorganic nature than
in organic nature. If then within inorganic nature we also distinguish the sphere
of mechanics from that of physics and chemistry in the narrower sense, the same
difference presents itself again. As is commonly acknowledged, mechanics is the
scientific discipline least capable of forgoing the assistance of mathematics; indeed
hardly a single step can be taken in mechanics without it, and for that reason
mechanics is also usually regarded, next to mathematics itself, as the exact sci-
ence par excellence. In this connection, though, it is necessary to recall again the
above comment concerning the coincidence of the materialist and the exclusively
mathematical standpoints. - IncidenraIly, after all that has been detailed here, it
must be called one of the most disruptive prejudices, precisely for knowledge of
an exact and thorough sort, if, as often happens, all difference and all determinacy
in the domain of the objective [des Gegenstiindlichen] are sought in what is merely
quantitative. To be sure, there is more to spirit than to nature: for instance, more
to the animal than to the plant. But one also knows very little about these objects
and their difference, if one merely stops short at this kind of a more or less and
does not proceed to construe them in the determinacy that is peculiar to them, a
determinacy that is here initially qualitative.
159
§ 100
Quantity, posited at first in its immediate relation to itself or in the deter-
mination of equality [Gleichheit] with itself as posited by attraction, is
continuous. According to the other determination contained in it, namely
that of the One, it is a discrete magnitude. The former quantity, however, is
equally discrete, since it is merely the continuity of the Many. The latter is
equally continuous, for its continuity is the One as the same in the many
Ones, i.e. the unity [of a mathematical unit].
Continuous and discrete magnitudes thus must not be regarded
as species, as though the determination of the one did not belong to
the other. Rather, they differ only in virtue of the fact that the same
whole is posited now under one and now under the other of its
determinations. 2. The antinomy of space, time, or matter {with
respect to their divisibility ad infinitum or their being composed of
1.160
The Encyclopedia Logic
indivisibles) is nothing but the assertion that quantity is now
continuous, now discrete. When space, time, etc. are posited only
with the determination of continuous quantity, they are divisible ad
infinitum; but with the determination of a discrete magnitude they
are in themselves divided and consist of indivisible ones. The one
is as one-sided as the other.
Addition. As the next result of being-for-itself, quantity contains the two sides
of its process, repulsion and attraction, as ideal moments within itself, and con-
sequently it is both continuous and discrete. Each of these twO moments equally
contains the other within itself, and hence there is neither a solely continuous nor
a solely discrete magnitude. If in spite of this one speaks of both as two particular,
mutually opposed species of magnitude, this is merely the result of our abstracting
reflection that in contemplating specific magnitudes ignores now the one and now
the other moment of the two contained in inseparable unity in the concept of
quantity. Thus it is said, for instance, that the space that this room occupies is
a continuous magnitude, and that these one hundred people who are gathered
together in it form a discrete magnitude. But space is continuous and discrete at
the same time, and we accordingly speak of points in space and then also divide
space; for instance, we divide a given extension into so many feet, inches, etc. This
can happen only under the supposition that space is in itself discrete as welL On
the other hand, the discrete magnitude consisting of a hundred people is simulta-
neously continuous as well, and what they have in common, namely the human
species that permeates all the individuals and connects them to each other, is that
in which the continuity of this magnitude is grounded.
b. Quantum
§ 101
Quantity, posited essentially with the exclusive determinacy that is con-
tained in it, is quantum, limited quantity.
Addition. Quantum is the existence of quantity, whereas pure quantity corre-
sponds to being and degree (to be considered shortly) to beingfor-itself. - As far
as the detail of the progression from pure quantity to quantum is concerned, this
progression is grounded in the fact that while in pure quantity the difference, as
a difference between continuity and discreteness, is at first on hand only in itself,
in the quantum, by contrast, this difference is posited, and indeed in such a way
that quantity now generally appears to be distinct or limited. As a result, however,
quantum simultaneously falls apart into an indeterminate assortment of quanta or
determinate magnitudes as welL Each of these determinate magnitudes, as distinct
from the others, forms a unity [the unity of a mathematical unit], while on the
The Encyclopedia Logic
16I
other hand, when considered in itself it is a Many. In this way, however, quantum
is determined as number.
§ 102
Quantum has its development and complete determinacy in number, which
contains the One as its element within itself and, as its qualitative moment,
the amount [Anz:ah~, which is the moment of discreteness, and the unity
[of a mathematical unit]' which is the moment of continuity.
In arithmetic, the kinds of calculation are usually listed as contingent
ways of treating numbers. If there is to be any necessity and thus
some rhyme and reason to them, it must lie in a principle, and that
principle can lie only in the determinations that are contained in
number itself This principle shall be briefly expounded here. - The
determinations of the concept of number are the amount (Anz:ahn
and the unity [of the mathematical unit], and number is the unity
of both. But the unity (of the mathematical unit], when applied
to empirical numbers, is merely the equality of them. Hence the
principle of the kinds of calculation must be to put numbers into
the relationship of the amount and the unity [of the mathematical
unit], and to produce the equality of these determinations.
Since the Ones or numbers are themselves indifferent towards
each other, the unity into which they are placed appears generally
to be an extraneous gathering together. For this reason, to calculate
generally means to count, and the difference between the kinds
of calculating resides exclusively in the qualitative make-up
[BeschaJfenheit] of the numbers that are being added together, and
the determination of the unity [of the mathematical unit]' and the
amount is the principle of their qualitative make-up.
To number or to generate number in general comes first, a matter
of taking arbitrarily many Ones together. - But calculation of a
particular sort is a matter of counting together items that are already
numbers, not the mere One.
Numbers are immediately and at first quite undetermined
numbers in general and, hence, unequal in general. Taking them
together or counting them is adding.
The next determination is that numbers are in general equal. Thus
they constitute a unity [Le. a mathematical unit] and there exists a
certain amount [Anz:ah~ of them. To count numbers such as these isThe Encyclopedia Logic
to multiply, in which case it does not matter how the determinations
of the amount and the unity [the mathematical unit] are distributed
to the two numbers or factors, i.e. which is taken to be the amount
and which the unity [the mathematical unit].
The third determinacy, finally, is the equality of amount and
unity [the mathematical unit]. Counting together the numbers
determined in this way is the raising of the power, and first of
all squaring. - The further raising of the power is the formal
continuation of the multiplication of number with itself, a
continuation that leads again to the indeterminate amount
[AnzaM. - Since perfect equality of the only difference on hand,
that of the amount and their unity, is attained in this third
determination, there cannot be more than these three kinds of
calculation. - To each of these ways of counting together there
corresponds the dissolution of numbers in accordance with the same
determinacies. Consequently, besides the three kinds listed, which
in that regard could be called positive, there also exist three negative
ones.
Addition. Because number in general is the quantum in its complete determi-
nacy, we use the quantum to determine not only so-called discrete magnitudes
but also so-called continuous magnitudes. For this reason, number must also be
utilized in geometry where the task is to indicate specific configurations of space
and their relationships.
c. Deg.ree
§ 103
The limit is identical to the whole of quantum itself. Insofar as it is, in
itself, manifold [vielfach], it is the extensive magnitude, but insofar as it is,
in itself, a simple determinateness, it is intensive magnitude, or deg.ree.
The difference between the continuous and discrete magnitudes and
the extensive and intensive ones consists in the fact that the former
apply to quantity in general, while the latter apply to the limit or
its determinacy as such. - Extensive and intensive magnitudes are
likewise not two species, each of which would contain a determinacy
that the other lacked. What is extensive magnitude is just as much
intensive magnitude, and vice versa.
The Encyclopedia Logic
Addition. Intensive magnitude or degree differs conceptually from extensive mag-
nitude or quantum. It must therefore be called illegitimate when, as often happens,
this distinction is not acknowledged and both forms of magnitude are identified
without further ado. This is notably the case in physics, for instance, when the
difference in specific weight is explained by saying that a body whose specific
weight is double that of another contains within the same space twice as many
material particles (atoms) as the other. The same goes for heat and light when the
different degrees of temperature and luminosity are supposed to be explained by a
greater or lesser amount of heat or light particles (or molecules). When confronted
with the illegitimacy of such explanations, physicists who avail themselves of them
tend to offer as an excuse that the in-itself of such phenomena (which as we know
is unknowable) is not supposed to be decided at all thereby, and that they avail
themselves of the expressions mentioned merely for the sake of their greater conve-
nience. But first, as far as the greater convenience is concerned, this is supposed to
be related to the easier application of the calculus. But it is by no means obvious
why intensive magnitudes, which find their determinate expression in number as
well, should not be as conveniently calculable as extensive magnitudes. It would
indeed be even more convenient to abandon all calculation as well as thinking
itself entirely. Further, against the excuse mentioned it must be remarked that by
engaging in these kinds of explanations one certainly goes beyond the domain of
perception and experience and enters the domain of metaphysics and speculation
(something that is otherwise declared to be idle or even pernicious). Experience
will indeed show that, if one of two purses filled with dollars is twice as heavy as
the other, then this is so because one of the purses contains two hundred and the
other only one hundred dollars. One can see these coins and altogether perceive
them with one's senses. By contrast, atoms, molecules, and things of this sort lie
outside the realm of sensory perception, and it is up to thinking to decide on
their admissibility and significance. But as mentioned earlier (§ 98, Addition), it
is the abstract understanding that fixes on the moment of the Many (contained
in the concept of being-for-itself) and does so in the shape of atoms, holding fast
to this moment as something ultimate. And it is the same abstract understanding
that in the case at hand again considers extensive magnitude to be the only form
of quantity, something that contradicts the untutored intuition just as much as
it does truly concrete thinking. Where intensive magnitudes present themselves
it does not acknowledge them in their characteristic determinacy, but instead,
relying on a hypothesis that is in itself untenable, tries to reduce them by force
to extensive magnitudes. Among the reproaches that have been brought against
recent philosophy, this one has been heard particularly frequently, namely that
it reduces everything to identity (and thus it has been mockingly called 'identity
philosophy'). From the discussion conducted here, it may be concluded that it
is precisely philosophy that presses for discriminating what is diverse in terms of
the concept as well as the experience of it, whereas it is professional empiricists
who elevate abstract identity to the highest principle of knowing and whose phi-
losophy for that reason should be more fittingly designated 'identity philosophy'.
Incidentally, it is quite right that just as little as there are magnitudes that are165
The Encyclopedia Logic The Encyclopedia Logic
exclusively continuous or discrete, so also there are no exclusively intensive or
extensive magnitudes, and that therefore the two determinations of quantity do
not stand opposite one another as self-standing species. Every intensive magnitude
is likewise extensive, and the same holds conversely. Thus, for instance, a certain
degree of temperature is an intensive magnitude to which as such an utterly simple
sensation also corresponds. If we then go and look at the thermometer, we find
how a certain expansion of the mercury column corresponds to this degree of
temperature and how that extensive magnitude changes in accordance with the
temperature as the intensive magnitude. It is the same with the domain of spirit: a
more intensive character reaches farther with its impact than a less intensive one. only one possibility (of being increased and decreased) and the insight into the
necessity of behaving in this way would be missing, quite apart from the fact that
in that case we would merely possess the representation and not the thought of
magnitude. By contrast, in the course of our logical development, not only has
quantity resulted as one stage of self-determining thinking, but it ~as also been
shown that it is inherent in the concept of quantity simply to propel Itself beyond
itself and that therefore we are here dealing not merely with something possible,
but with something necessary.
.
Addition 2. When reflective understanding is concerned with the infinite in
general, it tends to cling to the quantitative infinite progression above all. Now the
same thing that was mentioned earlier regarding the qualitative infinite progres-
sion holds good for this form of infinite progression as well, namely, that it is the
expression not of the true, but of the bad infinity that does not advance beyond
the mere ought and therefore in fact remains at a standstill in the finite. More
specifically, as far as the quantitative form of this finite progression is concerned -
something 5pinoza rightly designates a merely imagined infinite (infinitum imag-
inationis) - even poets (notably Haller and Klopstock) have frequently made use
of this representation in order to illustrate by means of it not only the infinity
of nature but also that of God himself. In Haller, for instance, we find a famous
description of God's infinity, which reads:
§ 104
In the [concept of] degree, the concept of quantum is posited. It is the
magnitude as indifferently for itself and simple, but in such a way that
it has the determinacy through which it is quantum entirely outside itself
in other magnitudes. With this contradiction, namely that the indifferent
limit that is for itself is the absolute externality, the infinite quantitative
progression is posited, - an immediacy that immediately changes over into
its opposite, i.e. into being mediated (i.e. transcending the quantum just
posited) and vice versa.
Number is thought, but thought as a being that is utterly external
to itself. It does not belong to intuition because it is thought, but
it is the thought that has the externality of intuition for its
determination. - For this reason, not only can quantum be increased
or decreased to infinity, it is through its concept this propelling
[Hinausschicken] of itself beyond itself. The infinite quantitative
progression is likewise the thoughtless repetition of the same
contradiction that quantum is in general and, when quantum is
posited in its determinacy, of the same contradiction that degree is.
Regarding the redundancy of expressing this contradiction in the
form of an infinite progression, Zeno rightly says in Aristotle: 'it is
the same thing to say something once and to be saying it always.'21
Addition I. When in mathematics, following the usual definition mentioned
earlier (§ 99), magnitude is designated as what can be increased or decreased -
and there is nothing objectionable about the underlying intuition here - the
question nevertheless still remains of how we come to assume something that is
capable of being increased and decreased. If one were to appeal simply to experience
to answer this question, this would not be sufficient, since it might prove to be
11
Translawrs' note: See Hermann Diels and Walther Kranz, Die Fragmente der Vorsokratiker 29
(Zenon) B1, Weidmann'sche Verlagsbuchhandlung, 1992.
I amass colossal numbers,
Millions of mountains,
I pile time upon time
And world upon worlds galore,
And when from this terrifying height
With vertigo I look to you again:
If all the mighty numbers
Were increased thousand-fold,
They would not even be a part of you.
So here we have, first, that constant propelling of quantity and, more specifi-
cally, of number beyond itself that Kant describes as horrifying, although the
actual horror is perhaps only the boredom of a limit constantly being set and
then sublated so that, as a result, one does not make any headway. But, further-
more, the poet mentioned aptly adds to this description of the bad infinity the
conclusion:
I remove them, and you lie before me entire
_ thereby making it explicit that the true infinite is not to be regarded as something
merely beyond the finite and that, to attain consciousness of it, we must renounce
that progressus in infinitum.
Addition 3. As is well known, Pythagoras philosophized in numbers and took
number to be the fundamental determination of things. At first glance, this way
of construing things necessarily appears quite paradoxical, even crazy to ordinary
consciousness, and the question therefore arises what is to be made of it. To
answer this question it must be remembered first that the task of philosophyThe Encyclopedia Logic The Encyclopedia Logic
generally consists in tracing things back to thoughts, and specifically to determinate
thoughts. Now number is, of course, a thought and, indeed, the thought that stands
~loses: to the sensory [sphere] or, put more precisely, the thought of the sensory
Itself, Insofar as we understand by it in general the [way things are] outside of one
another and a multiplicity. In the attempt to construe the universe as number we
thus recognize the first step towards metaphysics. Pythagoras stands in the history
of p~ilosophy, as is known, midway between the Ionian philosophers and the
Eleatlcs. Now while the former did not budge, as Aristotle has already remarked,
from regarding the essence of things as something material (as a hyle), and the latter,
Parmenides in particular, advanced to pure thinking in the form of being, it is the
Pythagorean philosophy whose principle forms the bridge, so to speak, between
the sensory and the supersensory. From this it follows what one should think
about the view of those who believe that Pythagoras obviously went too for since
~e construed th~ essence of things as mere numbers, and then noted that one may
Indeed count things (there being nothing objectionable about that) but that things
are after all more than mere numbers. Now as far as the' more is concerned ~ that
is attributed to things, it may be readily admitted that things are more than mere
n~m~ers, but w~at matters. is what is to be understood by this 'more'. In keeping
with Its standpOint, the ordinary sensory consciousness will not hesitate to answer
the que.stion raised ~ere by pointing to sensible perception and to note accordingly
that things are not Just countable but beyond that also visible, capable of being
smelt, felt, etc. .Expressed in our modern terms, the reproach made against the
Pythagorean philosophy would thus be reducible to saying that it is too idealistic.
But it in fact behaves in exactly the opposite way, as can already be gleaned from
wha: ~as noted earlier.about the historical position of the Pythagorean philosophy.
For If It must be admitted that things are more than mere numbers, this is to be
understood in such a way that the mere thought of number does not yet suffice to
express the specific essence or the concept of things. Instead of claiming, therefore,
that Pythagoras went too for with his philosophy of numbers one should say, on
the contrary, that he did not go for enough, and the Eleatics had already made
the next step to~ards pure thinking. - Furthermore, there are, if not things, at
least s~ates of things and natural phenomena in general whose determinacy rests
essentl~ly on specific numbers and proportions of numbers. This is notably the
case With the difference between sounds and their harmonic accord. The familiar
~tory is told that the perception of this phenomenon motivated Pythagoras to
Interpret the essence of things as number. Now although in the interest of science
it is crucial to trace those appearances that are based on specific numbers back to
them, it is in no way legitimate to regard the determinacy of thought generally
as a merely numerical determinacy. To be sure, one may be induced initially to
connect the most general thought-determinations to the first numbers and thus to
say that one is the simple and immediate, two the difference and mediation, and
~hree t?e unity ?f both. These c?mbinations are, however, quite external, and it
IS not I?herent In the numbers named, as such, to be the expression of just these
determinate thoughts. Moreover, the further one progresses in this manner, the
more the sheer arbitrariness of combining specific numbers with specific thoughts becomes apparent. Thus, for instance, one might regard 4 as the unity of I and 3
and the thoughts connected with them. But 4 is just as much the doubling of 2,
and in the same vein 9 is not only the square of 3 but also the sum of 8 and I, of
7 and 2, and so on. If even today certain secret societies place such a great weight
on all kinds of numbers and figures, this is to be regarded on the one hand as a
harmless game, but on the other hand as a sign of ineptness in thinking. True, one
likes to say that a deeper meaning lies hidden behind such things and that one can
think much to oneself thereby. However, what matters in philosophy is not that
one can think something, but that one actually thinks, and the true element of
thought must be sought not in arbitrarily chosen symbols but only in the thinking
itself.
166
§ 105
Quantum's being external to itself in the determinacy of its being-Jor-itself
constitutes its quality. In being external, it is precisely itself and related
to itself. The externality, i.e. the quantitative, and the being-for-itself, i.e.
the qualitative, are united therein. - Posited thus in itself[an ihm selbst],
quantum is quantitative proportion [Verhaltnis] - a determinacy that is just
as much an immediate quantum - i.e. the exponent as it is mediation,
i.e. the relation [Beziehung] of a given quantum to another, these being the
two sides of the proportion that at the same time are not to be taken in
their immediate value, but whose value lies exclusively in this relation.
Addition. The quantitative infinite progression initially appears to be an inces-
sant propelling of number beyond itself. Looked at more closely, however, the
quantity proves to be recursive [zu sich selbst zUrUckkehrend] in this progression,
for what is contained therein in terms of thought is that number is determined
generally by number, and this yields the quantitative proportion. When we say,
for instance, 2 in relation to 4, we have two magnitudes that are not to be taken
to be valid in their immediacy as such, but for which the relation to each other
is alone at issue. This relation, however, i.e. the exponent of the proportion [i.e.
2 in this case] is itself a magnitude that distinguishes itself from the magnitudes
that are related to each other in that the proportion changes when they change.
By contrast, the proportion is indifferent to the change of both of its sides and
remains the same as long as the exponent does not change. This is why instead
of 2:4 we can also put 3:6 in its place without altering the proportion, since the
exponent '2' remains the same in both cases.
§ 106
The sides of the proportion are still immediate quanta and the qualitative
and quantitative determinations are still external to each other. But as
for what they truly are, that the quantitative in its externality is itself168
The Encyclopedia Logic The Encyclopedia Logic
the relation to itself, or that being-for-itself and the indifference of the
determinacy are united, this is measure.
Addition. Due to the dialectical movement through its moments considered so far, quantity has proven to be the return to quality. The concept of quantity
was initially that of sublated quality, i.e. the determinacy that was not identical
with being, but instead indifferent, merely external to it. It is this concept that
also underlies mathematics' customary definition of magnitude, as mentioned
earlie~,. n~ely to be what can ~e increased and decreased. Now, according to this
definltJon It can seem at first as If magnitude were merely that which is alterable in
g:neral (for to in~rease ~ well as to decrease just means to determine a magnitude
differently). But If that IS the case, then magnitude would not be distinct from
existence (i.e. the second stage of quality) since according to its concept it is equally
alterable and the content of that definition would then have to be made complete
in such a way that in quantity we have something alterable that despite its alteration
remains the same. The concept of quantity thus proves to contain a contradiction
within itself, and it is this contradiction that constitutes the dialectic of quantity.
But the result of this dialectic is by no means a mere return to quality, as if the
latter was the true and quantit:f' by contrast the untrue. The result is instead'the
unity and truth of these two, qualitative quantity, or measure. - In this context,
it should be noted that when we are operating with quantitative determinations
in our examination of the objective [gegemtandlich] world it is indeed always
already measure that we have in view as the goal of such an examination. This is
indicated, moreover, in our language by the fact that the process of ascertaining
quantit~tive determi~ations and relationships is something that we designate as
measunng. Thus, for illS rance, one measures the lengths of various strings that are
made to vibrate with a view to the qualitative difference of the sounds produced
br t?e vib~ation, i~sofar as that d!fference corresponds to the difference in length.
SImilarly, III chemistry the quantity of substances [Stoffi] that are combined with
one another is ascertained in order to come to know the measurements that
condition these combinations, i.e. those quantities that underlie specific qualities.
In statistics, too, one deals with numbers but they are of interest only because of
the qualitative results conditioned by them, By contrast, the mere ascertaining of
numbers as such (without the guiding perspective specified here) rightly counts as
~n empty curiosity that is unable to satisfY either any theoretical or any practical
mterest. Addition. Measure as the unity of quality and quantity is thus at the same time
completed being. When we speak of being, it appears at first to be entirely abstract
and without determination. But being is essentially this, namely to determine
itself, and it attains its complete determinacy in measure. Measure can also be
regarded as a definition of the absolute, and it has been said accordingly that
God is the measure of all things. It is also this intuition that forms the basic tone
of so many ancient Hebraic songs in which the glorification of God essentially
amounts to the fact that he sets the limit for all things, for the sea and the land, the
rivers and mountains, and likewise for the various species of plants and animals. -
In the religious consciousness of the Greeks we find the divinity of measure in
closer connection to the ethical and represented as nemesis. This representation
contains the notion in general that everything human - wealth, honour, power,
and likewise joy, pain, etc. has its specific measure that, if overstepped, leads
to ruin and demise. - Now as far as the occurrence of measure in the objec-
tive [gegemtandlich] world is concerned, we find in nature initially concrete exis-
tences [Existenzen] whose essential content is formed by measure. This is notably
the case with the solar system that we are generally to regard as the realm of
free measures. Proceeding further in the contemplation of inorganic nature, mea-
sure recedes into the background, so to speak, insofar as the qualitative and quan-
titative determinations on hand here prove to be indifferent towards each other in
multiple ways. Thus, for instance, the quality of a rock or a river is not tied to a
specific magnitude. Looked at more closely, however, we find that even objects like
those mentioned are not absolutely measureless, for in a chemical analysis the water
in a river and the several components of a rock turn out to be qualities that are
again conditioned by quantitative proportions of the substances [Stoffi] contained
in them. But measure becomes decidedly more prominent again for immediate
intuition with organic nature. The various genera of plants and animals possess a
certain measure both overall and in their individual parts. Note should be taken,
moreover, of the circumstance that the less perfect organic forms that are closer
to inorganic nature distinguish themselves from the higher ones in part by the
greater indeterminacy of their measure. Thus, for example, among the petrifacts
we find so-called Ammon horns that are recognizable only in a microscope and
others that can be as large as a carriage wheel. The same indeterminacy of measure
also shows itself in many plants that are at a low level of organic development, as
is the case, for instance, with ferns.
C. MEASURE § 108
§ 107
Measure is the qualitative quantum, at first in the immediate sense as a
quantum with which an existence or a quality is bound up.
22
Translators' note: Taking Qualitiit here to be a ~iscue and that Quantitiit is meant.
Insofar as quality and quantity only exist in immediate unity in measure,
their difference emerges in them in an equally immediate manner. Insofar
as this is the case, the specific quantum is in part mere quantum, and
existence is capable of increase and decrease without the measure being
sublated thereby (the measure is in this respect a rule). In part, however,
the alteration of the quantum is likewise an alteration of quality.The Encyclopedia Logic The Encyclopedia Logic
Addition. The identity of quality and quantity present in measure is initially only
in itselfbut not yet posited. On account of this, these two determinations whose unchanged anymore. The constitution of a small Swiss canton does not fi~ a great
empire, and equally unsuitable was the constitution of the Roman Repubhc when
transferred to small German cities of the [Holy Roman] Empire.
170
unity is measure can render them each valid for itself such that, on the one hand,
the quantitative determinations of existence can be altered without its quality
being thereby affected but that, on the other hand, this process of indiscriminately
increasing and decreasing has its limits and the quality is altered by overstepping
those limits. Thus, for instance, the degree of the temperature of water is initially
irrelevant to its drop-forming fluidity. But in the course ofincreasing or decreasing
the temperature of water in its fluid, drop-forming state, a point is reached where
this state of cohesiveness changes qualitatively and water is transformed into vapour
on the one hand and ice on the other. When a quantitative change occurs, it appears
at first as something quite innocuous, and yet there is still something else hidden
behind it and this seemingly innocuous alteration of the quantitative is so to speak
a ruse (List] through which the qualitative is captured. The antinomy of measure
contained herein is something that the Greeks have already illustrated in many
different guises, such as, for example, in the question of whether one grain of wheat
makes a heap of wheat or in that other question of whether plucking one hair from
the tail of a horse makes a bald tail. While one may be initially inclined to respond
negatively to these questions in view of the nature of quantity as an indifferent
and external determinacy of being, one will nonetheless quickly have to admit that
such indifferent increasing and decreasing also has its limits and that there finally
comes a point where the result of the continued addition of only one grain at a
time is a heap of wheat and the result of the continued plucking of one hair at a
time is a bald tail. Similar to these examples is that tale of a peasant who continued
to increase by one ounce at a time the load of his donkey that was trotting along
cheerfully until it eventually collapsed under what had become an unbearable
burden. One would be very wrong if one were to declare such things to be merely
idle grammar school chatter, for they have in fact to do with thoughts with which
it is of great importance to be familiar also in a practical and more specifically in
an ethical connection. Thus, for instance, in connection with the expenditures we
make there is at first a certain latitude within which a more or less does not matter.
If, however, a certain measure, which is determined by the respective individual
situation, is overstepped, the qualitative nature of measure makes itself felt (in the
same way as in the previously mentioned example of the different temperatures
of water); and what just now had to be regarded as good budgeting turns into
stinginess or squandering. - The same point finds its application to politics as
well and, indeed, in the sense that the constitution of a state must be regarded as
both independent of and dependent on the extent of its territory, the number of
its inhabitants, and other such quantitative determinations. For instance, looking
at a state with a territory of one thousand square miles and a population of four
million inhabitants one will at first have to admit without hesitation that a few
square miles of territory or a few thousand inhabitants more or less cannot have a
significant influence on the constitution of such a state. But in contrast to this, it
is no less unmistakable that in the continued enlargement or shrinking of a state
there finally comes a point when because of this quantitative alteration, apart from
all other circumstances, the qualitative (aspect] of the constitution cannot remain
17 1
§ 109
The measureless is initially this process of a measure, by virtue of its quan-
titative nature, of passing beyond its qualitative determinacy. However,
since the other quantitative relationship, [namely,] the measurelessness of
the first, is equally qualitative, the measureless is likewise a measure. Both
transitions, from quality to quantum and from the latter to the former, can
be represented again as an infinite progression - as the suspending fAufheben]
and re-establishing of the measure in the measureless.
Addition. As we have seen, quantity is not only capable of alteration, i.e. of
increase and decrease; instead it is as such the process of stepping out (das Hinaus-
schreiten] beyond itself in general. Quantity confirms this its natute - in measure
as well. But as the quantity on hand in measure oversteps a certain limit, the qual-
ity corresponding to it is thereby likewise sub!ated. :x'ith th~s, however, i~ ~s n~t
quality in general that is negated, but only thIS speCIfic quality whose posmon IS
immediately re-occupied by another quality. This process of the measure proves to
be alternately the mere alteration of quantity and then a tipping over (Umschlagen]
of quantity into quality as well. One can make this process intuitively clear by using
the image of a knotted line. We first find knotted lines such as these in nature under
many different forms. Earlier we considered the qualitatively different states of the
aggregation of water that are dependent on increase and decrease. It is similar with
the different stages of the oxidation of metals. The differences between sounds
can also be regarded as an example of transforming what is at first merely quan-
titative into a qualitative alteration, a transforming that occurs in the process of
measure.
§ 110
What in fact happens here is that the immediacy, which still belongs to
the measure as such, is sublated. Quality and quantity themselves are first
in the measure as immediate, and it is merely their re14tive identity. Yet
measure turns out to sublate itself in the measureless. The latter, while it
is the negation of measure, is nonetheless itself the unity of quantity and
quality, and hence displays itself just as much as simply coming together
with itself.
§ III
The infinite, the affirmation as negation of negation, now has for its
sides quality and quantity instead of the more abstract sides of being andThe Encyclopedia Logic
nothing, something and an other, and so on. Those [sides of the infinite]
have made the transition (ex) first of quality into quantity (§ 98) and of
quantity into quality (§ 105) and been demonstrated by this means to be
negations. (~) But in their unity (i.e. measure) they are at first distinct and
one is only by being mediated by the other. And (y), after the immediacy of
this unity has proved to be self-sublating, this unity is now posited as what
it is in itself, as simple relation-to-itself that contains being in general and
its sublated forms within itself. - Being or the immediacy that is mediation
with itself and relation to itself through the negation of itself, and thus
equally mediation that sublates itself towards relation to itself, i.e. towards
immediacy, is essence.
